\section{Conclusion}\label{sec:conclusion}

We reproduced, audited and extended Qiao \emph{et al.}'s certificate-based aggregate signature scheme for
IIoT, with every result produced by running code on two independent elliptic-curve libraries. Three
conclusions stand out.

\begin{enumerate}
\item \textbf{Verma's scheme is broken more simply than the paper says.} Beyond the paper's
      malicious-KGC attack, which we reproduced, anyone holding a public key can forge, contradicting the
      paper's Table~II. The underlying rule, that every commitment in the verification equation must be
      hashed, explains both attacks and both of the repair's hash bindings.
\item \textbf{The repair works, but costs more than printed.} Its own verification equation costs about
      twice the row the paper prints, and about $2.4$ times Verma's in measured time. A best-effort
      optimisation of both schemes does not change that.
\item \textbf{Deploying the repair still needs care.} Pre-verify signatures at the aggregator, derive
      nonces with a hedged construction, and add verifier-checked freshness; the paper's replay-resistance
      and forward-security claims do not hold.
\end{enumerate}

\paragraph{Future work.} Three directions remain open: an audit of the paper's security proof,
measurements on constrained sensor hardware, and a comparison with Verma \emph{et al.}'s original
specification.
